% answers.tex -- answers and worked solutions for booklet.tex (Module 8 MC booklet), same
% layout (quiz.sty). One practiceans per quizq, in the same order (70).
% Every number: checks/m8-mc-booklet/numbers.py
% Build: TEXINPUTS=.: xelatex answers.tex (twice)
\documentclass[12pt]{article}
\usepackage{mathtools}
\usepackage{quiz}
\hypersetup{pdftitle={Multiple Choice Booklet - Module 8 - From the Universe to the Atom - Answers}}
% links between the two PDF files are not possible, so question headings are plain text here
\renewcommand\hyperlink[2]{#2}
\newcommand\anstopic[1]{\par\Needspace{8\baselineskip}\vspace{4mm}%
  \begin{center}{\headfont\bfseries\fontsize{14}{18}\selectfont #1\par}\end{center}\vspace{2pt}}
\newcommand\nuc[3]{\ensuremath{\prescript{#1}{#2}{\mathrm{#3}}}}

\begin{document}
\quiztitle{Year 12 Physics}{Multiple Choice Booklet - Module 8 - From the Universe to the Atom: Answers}

\begin{center}
  \renewcommand\arraystretch{1.45}
  \begin{tabular}{|l|*{14}{>{\centering\arraybackslash}p{6.2mm}|}}
    \hline
    \textbf{Q} & 1 & 2 & 3 & 4 & 5 & 6 & 7 & 8 & 9 & 10 & 11 & 12 & 13 & 14\\ \hline
    \textbf{Ans} & C & A & A & B & C & D & A & B & C & D & B & D & A & C\\ \hline\hline
    \textbf{Q} & 15 & 16 & 17 & 18 & 19 & 20 & 21 & 22 & 23 & 24 & 25 & 26 & 27 & 28\\ \hline
    \textbf{Ans} & D & C & B & B & A & D & A & C & B & B & A & C & D & D\\ \hline\hline
    \textbf{Q} & 29 & 30 & 31 & 32 & 33 & 34 & 35 & 36 & 37 & 38 & 39 & 40 & 41 & 42\\ \hline
    \textbf{Ans} & B & A & C & D & B & C & B & A & D & A & D & C & A & B\\ \hline\hline
    \textbf{Q} & 43 & 44 & 45 & 46 & 47 & 48 & 49 & 50 & 51 & 52 & 53 & 54 & 55 & 56\\ \hline
    \textbf{Ans} & C & D & A & B & C & D & B & A & A & D & C & B & C & A\\ \hline\hline
    \textbf{Q} & 57 & 58 & 59 & 60 & 61 & 62 & 63 & 64 & 65 & 66 & 67 & 68 & 69 & 70\\ \hline
    \textbf{Ans} & D & B & A & A & B & D & D & C & C & B & D & A & B & C\\ \hline
  \end{tabular}
\end{center}
Topic 1: Questions 1--14. \quad Topic 2: 15--28. \quad Topic 3: 29--42. \quad Topic 4: 43--56. \quad
Topic 5: 57--70.

Useful relationships: $v = H_0 d$, \ $\dfrac{1}{\lambda} = R\left(\dfrac{1}{n_f^2} - \dfrac{1}{n_i^2}\right)$,
\ $\lambda = \dfrac{h}{mv}$, \ $N = N_0 e^{-\lambda t}$, \ $\lambda = \dfrac{\ln 2}{t_{1/2}}$, \ $E = mc^2$,
\ $\qty{1}{u} = \qty{931.5}{MeV}/c^2$.

\begin{practiceanswers}

% =====================================================================
\anstopic{Topic 1 - Origins of the Elements}

% ---- Q1
\begin{practiceans}
  \correct{C}
  As the universe expanded and cooled, energy turned into matter. First came quarks (and electrons).
  After about a millionth of a second, quarks joined to form protons and neutrons. In the first few
  minutes, protons and neutrons fused into light nuclei (deuterium, helium and a little lithium). Neutral atoms formed only after about
  \num{380000} years, when it was cool enough for nuclei to hold on to electrons.
  \begin{itemize}
    \item \xmark\ A: this is the reverse order. \quad \xmark\ B and D: each puts a particle before its
          own building blocks.
  \end{itemize}
\end{practiceans}

% ---- Q2
\Needspace{9\baselineskip}
\begin{practiceans}
  \correct{A}
  The cosmic microwave background (CMB) is radiation left over from the hot early universe, released
  when atoms formed. Expansion has stretched it to microwave wavelengths (a \qty{2.7}{K} black body). The
  Big Bang theory predicted it, and it comes almost equally from all directions.
  \begin{itemize}
    \item \xmark\ B, C and D: these are true, but they are about stars and planets, not the origin of
          the universe.
  \end{itemize}
\end{practiceans}

% ---- Q3
\Needspace{9\baselineskip}
\begin{practiceans}
  \correct{A}
  Hubble found that the recession velocity is proportional to distance, $v = H_0 d$: a straight line
  through the origin. More distant galaxies move away faster, so the universe is expanding.
  \begin{itemize}
    \item \xmark\ B: the velocity is not the same for all galaxies.
    \item \xmark\ C and D: the velocity increases with distance; it does not decrease.
  \end{itemize}
\end{practiceans}

% ---- Q4
\Needspace{9\baselineskip}
\begin{practiceans}
  \correct{B}
  $H_0 = \dfrac{v}{d} = \dfrac{8400}{120} = \qty{70}{km.s^{-1}.Mpc^{-1}}$.
  \begin{itemize}
    \item \xmark\ A: this is $\frac{d}{v}$ (upside down). \quad \xmark\ C: this is just the velocity.
    \item \xmark\ D: this is $v \times d$.
  \end{itemize}
\end{practiceans}

% ---- Q5
\Needspace{9\baselineskip}
\begin{practiceans}
  \correct{C}
  Convert $H_0$ to SI units: $H_0 = \dfrac{\qty{70e3}{m.s^{-1}}}{\qty{3.09e22}{m}} = \qty{2.27e-18}{s^{-1}}$.
  \[ t = \frac{1}{H_0} = \qty{4.4e17}{s} = \frac{\num{4.4e17}}{\num{3.16e7}} = \qty{1.4e10}{years} \]
  \begin{itemize}
    \item \xmark\ A: this is the number of seconds, not years.
    \item \xmark\ B: did not change km into m. \quad \xmark\ D: wrong power of ten.
  \end{itemize}
\end{practiceans}

% ---- Q6
\Needspace{9\baselineskip}
\begin{practiceans}
  \correct{D}
  A red giant is cool (red) but very luminous because it is huge: top right of the diagram (W). A white
  dwarf is hot but faint because it is tiny: bottom left (X). Y is a hot, massive main-sequence star and
  Z is a cool, faint main-sequence star (a red dwarf).
  \begin{itemize}
    \item \xmark\ A, B and C: Y and Z are on the main sequence, and X is not a red giant.
  \end{itemize}
\end{practiceans}

% ---- Q7
\Needspace{9\baselineskip}
\begin{practiceans}
  \correct{A}
  The most massive (hottest, bluest) stars use up their fuel first and leave the main sequence first.
  As a cluster ages, the turn-off point moves down the main sequence to cooler, less massive stars.
  Cluster P still has O and B stars, so it is youngest; R has lost all stars hotter than G stars, so it
  is oldest.
  \begin{itemize}
    \item \xmark\ B: this is oldest to youngest. \quad \xmark\ C and D: do not follow the order of the
          turn-off points.
  \end{itemize}
\end{practiceans}

% ---- Q8
\Needspace{9\baselineskip}
\begin{practiceans}
  \correct{B}
  The stars in a cluster formed from the same cloud of gas at about the same time, so they have about
  the same age (and the same distance from Earth and the same starting composition). They have different
  masses, so they lie at different places on the diagram.
  \begin{itemize}
    \item \xmark\ A, C and D: the stars have a range of masses, luminosities and temperatures.
  \end{itemize}
\end{practiceans}

% ---- Q9
\Needspace{9\baselineskip}
\begin{practiceans}
  \correct{C}
  The same spectral class means the same surface temperature. A star with the same temperature but much
  greater luminosity must have a much larger surface area, so a larger radius (for example, a giant
  compared with a main-sequence star).
  \begin{itemize}
    \item \xmark\ A: the same spectral class means the same temperature.
    \item \xmark\ B: distance changes how bright a star \emph{looks}, not its luminosity.
    \item \xmark\ D: age does not explain the difference.
  \end{itemize}
\end{practiceans}

% ---- Q10
\Needspace{9\baselineskip}
\begin{practiceans}
  \correct{D}
  $\nuc{2}{1}{H} + \nuc{1}{1}{H} \rightarrow \nuc{3}{2}{He} + \gamma$. The mass numbers add to 3 and the
  atomic numbers add to 2, so the product is helium-3.
  \begin{itemize}
    \item \xmark\ A: helium-4 is made in the last step, when two helium-3 nuclei fuse.
    \item \xmark\ B: the atomic number must be $1 + 1 = 2$. \quad \xmark\ C: the mass number must be
          $2 + 1 = 3$.
  \end{itemize}
\end{practiceans}

% ---- Q11
\Needspace{9\baselineskip}
\begin{practiceans}
  \correct{B}
  In the CNO cycle, carbon, nitrogen and oxygen nuclei act as catalysts for turning hydrogen into helium.
  Because these nuclei have larger charges, the cycle needs higher core temperatures, so it dominates in
  main-sequence stars more than about 1.3 times the mass of the Sun. Stars like the Sun and smaller use
  mainly the proton--proton chain.
  \begin{itemize}
    \item \xmark\ A and C: smaller stars use the proton--proton chain.
    \item \xmark\ D: white dwarfs have no fusion in their cores.
  \end{itemize}
\end{practiceans}

% ---- Q12
\Needspace{9\baselineskip}
\begin{practiceans}
  \correct{D}
  $E = 0.0287 \times 931.5 = \qty{26.7}{MeV}$. (This uses atomic masses, so it includes the energy
  released when the positrons annihilate with electrons.)
  \begin{itemize}
    \item \xmark\ A: divided by $c^2$. \quad \xmark\ B: this is the energy in joules
          (\qty{4.3e-12}{J}), labelled as MeV.
    \item \xmark\ C: did not convert the mass defect from u to MeV.
  \end{itemize}
\end{practiceans}

% ---- Q13
\Needspace{9\baselineskip}
\begin{practiceans}
  \correct{A}
  A more massive main-sequence star has a hotter core, so fusion is much faster: it is far more
  luminous. It uses up its hydrogen much faster than it gains extra fuel from its larger mass, so it
  spends a much shorter time on the main sequence.
  \begin{itemize}
    \item \xmark\ B, C and D: each has at least one part the wrong way round.
  \end{itemize}
\end{practiceans}

% ---- Q14
\Needspace{9\baselineskip}
\begin{practiceans}
  \correct{C}
  Iron-56 has about the highest binding energy per nucleon, so making heavier nuclei by fusion
  \emph{absorbs} energy. Heavier elements are made mainly by neutron capture (followed by beta decay):
  slowly in red giant stars, and rapidly in supernovae and neutron-star mergers.
  \begin{itemize}
    \item \xmark\ A: the Sun's core fuses hydrogen into helium.
    \item \xmark\ B: fusion beyond iron does not release energy.
    \item \xmark\ D: the Big Bang made mostly hydrogen and helium (and a little lithium).
  \end{itemize}
\end{practiceans}

% =====================================================================
\anstopic{Topic 2 - Structure of the Atom}

% ---- Q15
\begin{practiceans}
  \correct{D}
  Cathode rays were deflected by electric and magnetic fields in the direction of negative charge.
  Thomson measured their charge-to-mass ratio, which was more than 1000 times that of a hydrogen ion (the
  modern value is about 1840 times), and the same for every cathode material. So they are tiny negative particles found in all atoms: electrons.
  \begin{itemize}
    \item \xmark\ A and C: they are deflected by fields, so they are charged particles.
    \item \xmark\ B: they are deflected as negative charges.
  \end{itemize}
\end{practiceans}

% ---- Q16
\Needspace{9\baselineskip}
\begin{practiceans}
  \correct{C}
  Undeflected means $qE = qvB$, so $v = \dfrac{E}{B} = \dfrac{\num{2.0e4}}{\num{1.0e-3}} = \qty{2.0e7}{m.s^{-1}}$.
  \begin{itemize}
    \item \xmark\ A: this is $E \times B$. \quad \xmark\ B: used $B = \qty{1.0}{T}$.
    \item \xmark\ D: this is $\frac{E}{B^2}$.
  \end{itemize}
\end{practiceans}

% ---- Q17
\Needspace{9\baselineskip}
\begin{practiceans}
  \correct{B}
  $qvB = \dfrac{mv^2}{r}$, so $\dfrac{q}{m} = \dfrac{v}{Br} = \dfrac{\num{2.0e7}}{(\num{1.0e-3})(0.114)}
  = \qty{1.8e11}{C.kg^{-1}}$.
  \begin{itemize}
    \item \xmark\ A: used $v^2$. \quad \xmark\ C: this is $\frac{1}{Br}$ (forgot $v$).
    \item \xmark\ D: this is $\frac{m}{q}$ (upside down).
  \end{itemize}
\end{practiceans}

% ---- Q18
\Needspace{9\baselineskip}
\begin{practiceans}
  \correct{B}
  $E = \dfrac{V}{d} = \dfrac{450}{0.015} = \qty{3.0e4}{V.m^{-1}}$. Stationary: $qE = mg$, so
  \[ q = \frac{mg}{E} = \frac{(\num{4.9e-15})(9.8)}{\num{3.0e4}} = \qty{1.6e-18}{C}, \qquad
     n = \frac{q}{e} = \frac{\num{1.6e-18}}{\num{1.602e-19}} = 10 \]
  \begin{itemize}
    \item \xmark\ A: forgot to multiply by $g$. \quad \xmark\ C: wrong power of ten.
    \item \xmark\ D: used the plate spacing in cm instead of m.
  \end{itemize}
\end{practiceans}

% ---- Q19
\Needspace{9\baselineskip}
\begin{practiceans}
  \correct{A}
  Millikan found that the charge on every drop was a whole-number multiple of one smallest charge,
  $e = \qty{1.6e-19}{C}$. So charge is quantised, and $e$ is the charge on one electron.
  \begin{itemize}
    \item \xmark\ B: the mass came from combining $e$ with Thomson's $\frac{e}{m}$, not from the drops.
    \item \xmark\ C: drops could have either sign. \quad \xmark\ D: only multiples of $e$ were found.
  \end{itemize}
\end{practiceans}

% ---- Q20
\Needspace{9\baselineskip}
\begin{practiceans}
  \correct{D}
  In the plum-pudding model, the positive charge is spread through the whole atom, so its electric field
  is weak, and the electrons are far too light to deflect a heavy alpha particle much. So the alpha
  particles were expected to pass through with only tiny deflections.
  \begin{itemize}
    \item \xmark\ A and C: large deflections need a small, concentrated charge, which this model did not
          have.
    \item \xmark\ B: alpha particles were expected to pass through thin foil.
  \end{itemize}
\end{practiceans}

% ---- Q21
\Needspace{9\baselineskip}
\begin{practiceans}
  \correct{A}
  Only a very strong force could turn an alpha particle back. This needs the positive charge (and most of
  the mass) to be packed into a tiny nucleus. Because so few were turned back, the nucleus must be very
  small.
  \begin{itemize}
    \item \xmark\ B: this conclusion came from the \emph{many} alpha particles that passed straight
          through.
    \item \xmark\ C: this is Bohr's idea. \quad \xmark\ D: atoms are neutral.
  \end{itemize}
\end{practiceans}

% ---- Q22
\Needspace{9\baselineskip}
\begin{practiceans}
  \correct{C}
  The radiation was not deflected by fields (neutral) and was very penetrating. From the recoil of the
  protons (and of nitrogen nuclei in another test), Chadwick showed it was a neutral particle with about
  the same mass as a proton: the neutron.
  \begin{itemize}
    \item \xmark\ A: see Q23. \quad \xmark\ B: electrons are charged.
    \item \xmark\ D: alpha particles are charged and not very penetrating.
  \end{itemize}
\end{practiceans}

% ---- Q23
\Needspace{9\baselineskip}
\begin{practiceans}
  \correct{B}
  To knock protons out at the observed speeds, gamma photons would need far more energy (about
  \qty{50}{MeV}) than the reaction could supply. A particle with about the same mass as a proton can pass
  on its kinetic energy very efficiently in a collision, like one billiard ball hitting another.
  \begin{itemize}
    \item \xmark\ A and D: the radiation was neutral, so it was not deflected.
    \item \xmark\ C: nothing travels faster than light.
  \end{itemize}
\end{practiceans}

% ---- Q24
\Needspace{9\baselineskip}
\begin{practiceans}
  \correct{B}
  At speed $v$: $qE = qvB$. At $2v$, the magnetic force doubles to $2qE$, but the electric force stays $qE$
  (it does not depend on speed). They act in opposite directions, so the net force is $2qE - qE = qE$:
  \[ a = \frac{qE}{m} = \frac{(\num{1.602e-19})(\num{5.0e3})}{\num{1.673e-27}} = \qty{4.8e11}{m.s^{-2}} \]
  \begin{itemize}
    \item \xmark\ A: the forces balance only at speed $v$.
    \item \xmark\ C: this uses the magnetic force alone. \quad \xmark\ D: added the two forces.
  \end{itemize}
\end{practiceans}

% ---- Q25
\Needspace{9\baselineskip}
\begin{practiceans}
  \correct{A}
  Cathode rays are electrons (negative), so they are attracted towards the positive plate: upwards.
  \begin{itemize}
    \item \xmark\ B: this would be the direction for positive particles.
    \item \xmark\ C: charged particles are deflected by an electric field.
    \item \xmark\ D: the electric force is along the field lines (vertical here).
  \end{itemize}
\end{practiceans}

% ---- Q26
\Needspace{9\baselineskip}
\begin{practiceans}
  \correct{C}
  $N = A - Z = 238 - 92 = 146$.
  \begin{itemize}
    \item \xmark\ A: added $A$ and $Z$. \quad \xmark\ B: this is the mass number (protons + neutrons).
    \item \xmark\ D: this is the number of protons.
  \end{itemize}
\end{practiceans}

% ---- Q27
\Needspace{9\baselineskip}
\begin{practiceans}
  \correct{D}
  Thomson pictured the atom as a sphere of positive charge with the electrons embedded in it, like plums
  in a pudding.
  \begin{itemize}
    \item \xmark\ A: this is Rutherford's model. \quad \xmark\ B: this is Bohr's model.
    \item \xmark\ C: this is the quantum mechanical (Schrödinger) model.
  \end{itemize}
\end{practiceans}

% ---- Q28
\Needspace{9\baselineskip}
\begin{practiceans}
  \correct{D}
  At first, $qE = mg$. Doubling the potential difference doubles $E$, so the electric force becomes $2mg$
  upwards. The net force is $2mg - mg = mg$ upwards, so the drop accelerates upwards (at $g$, until air
  resistance becomes important).
  \begin{itemize}
    \item \xmark\ A: the forces are no longer balanced.
    \item \xmark\ B and C: the net force is upwards, not downwards.
  \end{itemize}
\end{practiceans}

% =====================================================================
\anstopic{Topic 3 - Quantum Mechanical Nature of the Atom}

% ---- Q29
\begin{practiceans}
  \correct{B}
  An orbiting electron is always accelerating (its direction keeps changing). In classical physics an
  accelerating charge radiates EM waves, so the electron would lose energy and spiral into the nucleus in
  a tiny fraction of a second. Real atoms are stable.
  \begin{itemize}
    \item \xmark\ A: the electric attraction holds the electrons; the size of the nucleus is not the
          problem.
    \item \xmark\ C: true (neutrons came later), but this does not affect stability.
    \item \xmark\ D: not a problem with the model.
  \end{itemize}
\end{practiceans}

% ---- Q30
\Needspace{9\baselineskip}
\begin{practiceans}
  \correct{A}
  Bohr's fixed energy levels explain why hydrogen emits only certain wavelengths: each line is a photon
  with $hf = E_i - E_f$. Rutherford's model predicted a continuous spread of emitted wavelengths as an
  electron spiralled inwards.
  \begin{itemize}
    \item \xmark\ B: Rutherford's model explained this. \quad \xmark\ C: both models explain this.
    \item \xmark\ D: this is about free electrons, not atomic models.
  \end{itemize}
\end{practiceans}

% ---- Q31
\Needspace{9\baselineskip}
\begin{practiceans}
  \correct{C}
  Bohr's postulates: (1) electrons exist only in certain stationary states, in which they do not radiate;
  (2) an electron moving between states emits or absorbs a photon with $hf = E_i - E_f$; (3) the angular
  momentum of an electron is quantised ($mvr = \frac{nh}{2\pi}$).
  \begin{itemize}
    \item \xmark\ A: this is the classical idea that Bohr rejected.
    \item \xmark\ B: neutrons were not discovered until 1932. \quad \xmark\ D: this is Thomson's model.
  \end{itemize}
\end{practiceans}

% ---- Q32
\Needspace{9\baselineskip}
\begin{practiceans}
  \correct{D}
  \[ \frac{1}{\lambda} = R\left(\frac{1}{2^2} - \frac{1}{3^2}\right) = (\num{1.097e7})(0.1389) = \qty{1.524e6}{m^{-1}},
     \qquad \lambda = \qty{6.56e-7}{m} = \qty{656}{nm} \]
  \begin{itemize}
    \item \xmark\ A: this is the $n = 2$ to $n = 1$ line. \quad \xmark\ B: added the two terms instead of
          subtracting.
    \item \xmark\ C: used $\frac{1}{n}$ instead of $\frac{1}{n^2}$.
  \end{itemize}
\end{practiceans}

% ---- Q33
\Needspace{9\baselineskip}
\begin{practiceans}
  \correct{B}
  The shortest wavelength means the largest photon energy ($E = \frac{hc}{\lambda}$), so the largest energy
  gap: $n = 2$ to $n = 1$ gives $-3.40 - (-13.6) = \qty{10.2}{eV}$.
  \begin{itemize}
    \item \xmark\ A: \qty{0.66}{eV}. \quad \xmark\ C: \qty{1.89}{eV}. \quad \xmark\ D: \qty{2.55}{eV}.
  \end{itemize}
\end{practiceans}

% ---- Q34
\Needspace{9\baselineskip}
\begin{practiceans}
  \correct{C}
  The possible transitions are $4 \to 3$, $4 \to 2$, $4 \to 1$, $3 \to 2$, $3 \to 1$ and $2 \to 1$: six
  different energy gaps, so six wavelengths.
  \begin{itemize}
    \item \xmark\ A: counts only the jumps that start at $n = 4$. \quad \xmark\ B: this is the number of
          levels.
    \item \xmark\ D: too many; only six pairs of levels exist.
  \end{itemize}
\end{practiceans}

% ---- Q35
\Needspace{9\baselineskip}
\begin{practiceans}
  \correct{B}
  $E = E_3 - E_1 = -1.51 - (-13.6) = \qty{12.1}{eV}$.
  \begin{itemize}
    \item \xmark\ A: this would ionise the atom. \quad \xmark\ C: this is the energy to reach $n = 2$.
    \item \xmark\ D: this is the size of the $n = 3$ energy level, not the gap.
  \end{itemize}
\end{practiceans}

% ---- Q36
\Needspace{9\baselineskip}
\begin{practiceans}
  \correct{A}
  $\lambda = \dfrac{h}{mv} = \dfrac{\num{6.626e-34}}{(\num{9.109e-31})(\num{2.0e6})} = \qty{3.6e-10}{m}$.
  \begin{itemize}
    \item \xmark\ B: used the mass of a proton. \quad \xmark\ C: this is the momentum $mv$.
    \item \xmark\ D: this is $h \times m \times v$.
  \end{itemize}
\end{practiceans}

% ---- Q37
\Needspace{9\baselineskip}
\begin{practiceans}
  \correct{D}
  $\lambda = \frac{h}{mv}$. At the same speed, the wavelength depends only on the mass. The neutron
  (\qty{1.675e-27}{kg}) is very slightly heavier than the proton (\qty{1.673e-27}{kg}), so its wavelength
  is very slightly (about 0.1\%) shorter.
  \begin{itemize}
    \item \xmark\ A: charge does not affect the de Broglie wavelength.
    \item \xmark\ B: the masses are slightly different. \quad \xmark\ C: the masses are almost equal.
  \end{itemize}
\end{practiceans}

% ---- Q38
\Needspace{9\baselineskip}
\begin{practiceans}
  \correct{A}
  $\lambda_\alpha = \dfrac{h}{(4m)v}$ and $\lambda_\mathrm{p} = \dfrac{h}{m(2v)}$, so
  $\dfrac{\lambda_\alpha}{\lambda_\mathrm{p}} = \dfrac{2mv}{4mv} = \dfrac{1}{2}$.
  The alpha particle has twice the momentum, so half the wavelength.
  \begin{itemize}
    \item \xmark\ B, C and D: compare the momenta: $4mv$ against $2mv$.
  \end{itemize}
\end{practiceans}

% ---- Q39
\Needspace{9\baselineskip}
\begin{practiceans}
  \correct{D}
  The shortest wavelength belongs to the largest momentum $mv$. Using $\lambda = \frac{h}{mv}$:
  A \qty{7.3e-10}{m}, B \qty{1.5e-10}{m}, C \qty{4.0e-10}{m}, D \qty{2.0e-11}{m}.
  \begin{itemize}
    \item \xmark\ A, B and C: these particles have less momentum than the proton in D.
  \end{itemize}
\end{practiceans}

% ---- Q40
\Needspace{9\baselineskip}
\begin{practiceans}
  \correct{C}
  Davisson and Germer fired electrons at a nickel crystal and found a diffraction pattern. The measured
  wavelength agreed with $\lambda = \frac{h}{mv}$, confirming de Broglie's idea that matter has wave
  properties.
  \begin{itemize}
    \item \xmark\ A: shown by Thomson. \quad \xmark\ B: shown by the photoelectric effect.
    \item \xmark\ D: shown by the Geiger--Marsden experiment.
  \end{itemize}
\end{practiceans}

% ---- Q41
\Needspace{9\baselineskip}
\begin{practiceans}
  \correct{A}
  De Broglie said an electron in an orbit is a standing wave. It is stable only if a whole number of
  wavelengths fits around the orbit: $2\pi r = n\lambda$. This gives exactly Bohr's condition
  $mvr = \frac{nh}{2\pi}$, explaining why only certain orbits are allowed.
  \begin{itemize}
    \item \xmark\ B: electrons cannot reach the speed of light.
    \item \xmark\ C: not part of the explanation. \quad \xmark\ D: electrons in stationary states do not
          radiate.
  \end{itemize}
\end{practiceans}

% ---- Q42
\Needspace{9\baselineskip}
\begin{practiceans}
  \correct{B}
  Bohr's model worked well for hydrogen (one electron), but it could not explain the spectra of atoms
  with more than one electron, the different brightness (intensity) of the lines, or the splitting of lines
  into close pairs (fine structure). It also mixed classical and quantum ideas without a reason for the
  stationary states.
  \begin{itemize}
    \item \xmark\ A: explaining the hydrogen spectrum was its great success.
    \item \xmark\ C: it included Rutherford's nucleus. \quad \xmark\ D: it made atoms stable.
  \end{itemize}
\end{practiceans}

% =====================================================================
\anstopic{Topic 4 - Properties of the Nucleus}

% ---- Q43
\begin{practiceans}
  \correct{C}
  Beta particles (electrons) are negative and very light, so they bend strongly towards the positive
  plate: P. Gamma rays are uncharged, so they go straight: Q. Alpha particles are positive but much
  heavier, so they bend only slightly towards the negative plate: R.
  \begin{itemize}
    \item \xmark\ A, B and D: each has at least one path wrong.
  \end{itemize}
\end{practiceans}

% ---- Q44
\Needspace{9\baselineskip}
\begin{practiceans}
  \correct{D}
  Work backwards. The nucleus before the alpha decay had $A = 7 + 4 = 11$ and $Z = 3 + 2 = 5$: boron-11.
  This was made by adding a neutron, so the original nucleus had $A = 11 - 1 = 10$, $Z = 5$: boron-10.
  \[ \nuc{10}{5}{B} + \nuc{1}{0}{n} \rightarrow \nuc{11}{5}{B} \rightarrow \nuc{7}{3}{Li} + \nuc{4}{2}{He} \]
  \begin{itemize}
    \item \xmark\ B: this is the nucleus formed \emph{after} the neutron is absorbed.
    \item \xmark\ A and C: the mass and atomic numbers do not balance.
  \end{itemize}
\end{practiceans}

% ---- Q45
\Needspace{9\baselineskip}
\begin{practiceans}
  \correct{A}
  Alpha decay removes 2 protons and 2 neutrons: $A$ falls by 4 and $Z$ falls by 2.
  $\nuc{226}{88}{Ra} \rightarrow \nuc{222}{86}{Rn} + \nuc{4}{2}{He}$.
  \begin{itemize}
    \item \xmark\ B: the mass number must fall by 4. \quad \xmark\ C: the atomic number must fall by 2.
    \item \xmark\ D: added the alpha particle instead of removing it.
  \end{itemize}
\end{practiceans}

% ---- Q46
\Needspace{9\baselineskip}
\begin{practiceans}
  \correct{B}
  In beta-minus decay, a neutron changes into a proton, emitting an electron and an antineutrino. The
  mass number stays the same and the atomic number increases by 1, so carbon-14 becomes nitrogen-14.
  \begin{itemize}
    \item \xmark\ A: this is beta-plus decay. \quad \xmark\ C: this is alpha decay.
    \item \xmark\ D: this is neutron emission.
  \end{itemize}
\end{practiceans}

% ---- Q47
\Needspace{9\baselineskip}
\begin{practiceans}
  \correct{C}
  X has $N = 144$, $Z = 90$; Y has $N = 143$, $Z = 91$. One neutron has become a proton ($N$ down 1,
  $Z$ up 1, $A$ unchanged): beta-minus decay.
  \begin{itemize}
    \item \xmark\ A: alpha decay lowers both $N$ and $Z$ by 2.
    \item \xmark\ B: beta-plus decay raises $N$ by 1 and lowers $Z$ by 1.
    \item \xmark\ D: gamma decay does not change $N$ or $Z$.
  \end{itemize}
\end{practiceans}

% ---- Q48
\Needspace{9\baselineskip}
\begin{practiceans}
  \correct{D}
  Each half-life halves the mass: $80 \rightarrow 40 \rightarrow 20 \rightarrow \qty{10}{g}$.
  \begin{itemize}
    \item \xmark\ A: after one half-life. \quad \xmark\ B: divided by 3. \quad \xmark\ C: after two
          half-lives.
  \end{itemize}
\end{practiceans}

% ---- Q49
\Needspace{9\baselineskip}
\begin{practiceans}
  \correct{B}
  After two half-lives, $\frac{1}{4}$ of the parent atoms remain, and the other $\frac{3}{4}$ have become
  daughter atoms. Parent : daughter $= \frac{1}{4} : \frac{3}{4} = 1 : 3$.
  \begin{itemize}
    \item \xmark\ A: compares the parent with the original number, not with the daughter.
    \item \xmark\ C: does not match any number of half-lives. \quad \xmark\ D: after one half-life.
  \end{itemize}
\end{practiceans}

% ---- Q50
\Needspace{9\baselineskip}
\begin{practiceans}
  \correct{A}
  From the graph, the mass halves from \qty{40}{g} to \qty{20}{g} in 4 days, so $t_{1/2} = 4.0$ days.
  $\lambda = \dfrac{\ln 2}{t_{1/2}} = \dfrac{0.693}{4.0} = \qty{0.17}{day^{-1}}$.
  \begin{itemize}
    \item \xmark\ B: this is $\frac{1}{t_{1/2}}$. \quad \xmark\ C: this is $t_{1/2} \times \ln 2$.
    \item \xmark\ D: this is the half-life.
  \end{itemize}
\end{practiceans}

% ---- Q51
\Needspace{9\baselineskip}
\begin{practiceans}
  \correct{A}
  The smallest decay constant means the longest half-life ($\lambda = \frac{\ln 2}{t_{1/2}}$). Graph A
  falls to $\frac{A_0}{2}$ at 3 hours; D at 2 hours; C at 1 hour.
  \begin{itemize}
    \item \xmark\ B: a straight line is not exponential decay, so it is not a valid decay curve (it also
          halves in only 2 hours).
    \item \xmark\ C and D: shorter half-lives, so larger decay constants.
  \end{itemize}
\end{practiceans}

% ---- Q52
\Needspace{9\baselineskip}
\begin{practiceans}
  \correct{D}
  $\lambda = \dfrac{\ln 2}{8.0} = \qty{0.0866}{day^{-1}}$, so
  $A = A_0 e^{-\lambda t} = 800\,e^{-0.0866 \times 10} = \qty{3.4e2}{Bq}$.
  \begin{itemize}
    \item \xmark\ A: this is the activity after 3 half-lives (24 days).
    \item \xmark\ B: used $e^{-t/t_{1/2}}$ (forgot the $\ln 2$). \quad \xmark\ C: assumed the activity
          falls in a straight line.
  \end{itemize}
\end{practiceans}

% ---- Q53
\Needspace{9\baselineskip}
\begin{practiceans}
  \correct{C}
  Mass defect: $2(1.007276) + 2(1.008665) - 4.001506 = \qty{0.030376}{u}$. \quad
  Binding energy $= 0.030376 \times 931.5 = \qty{28.3}{MeV}$.
  \begin{itemize}
    \item \xmark\ A: did not convert u to MeV. \quad \xmark\ B: this is the binding energy per nucleon.
    \item \xmark\ D: this is the rest energy of the whole nucleus.
  \end{itemize}
\end{practiceans}

% ---- Q54
\Needspace{9\baselineskip}
\begin{practiceans}
  \correct{B}
  The binding energy is the energy equivalent of the mass defect: $E_\mathrm{B} = \Delta m\,c^2$. So a
  larger binding energy always means a larger mass defect.
  \begin{itemize}
    \item \xmark\ A: usually true, but not always (for example, \nuc{56}{26}{Fe} has a larger binding energy
          than \nuc{56}{28}{Ni}, but a smaller mass).
    \item \xmark\ C: stability depends on the binding energy \emph{per nucleon}.
    \item \xmark\ D: not necessarily.
  \end{itemize}
\end{practiceans}

% ---- Q55
\Needspace{9\baselineskip}
\begin{practiceans}
  \correct{C}
  Energy released $=$ binding energy of the products $-$ binding energy of the reactants (a free neutron
  has no binding energy): $28.30 - (2.22 + 8.48) = \qty{17.6}{MeV}$.
  \begin{itemize}
    \item \xmark\ A: added all three. \quad \xmark\ B: the binding energy of helium-4 only.
    \item \xmark\ D: this is $8.48 - 2.22$.
  \end{itemize}
\end{practiceans}

% ---- Q56
\Needspace{9\baselineskip}
\begin{practiceans}
  \correct{A}
  Fission produces fast neutrons, but uranium-235 captures slow (thermal) neutrons much more easily. The
  moderator (water or graphite) slows the neutrons down by collisions, so the chain reaction can continue.
  \begin{itemize}
    \item \xmark\ B: this is the job of the control rods. \quad \xmark\ C: this is the coolant.
    \item \xmark\ D: this is the shielding.
  \end{itemize}
\end{practiceans}

% =====================================================================
\anstopic{Topic 5 - Deep Inside the Atom}

% ---- Q57
\begin{practiceans}
  \correct{D}
  A proton is a baryon made of three quarks (uud), so it is not fundamental.
  \begin{itemize}
    \item \xmark\ A and B: electrons and muons are leptons, which are fundamental.
    \item \xmark\ C: the photon is a fundamental force carrier (gauge boson).
  \end{itemize}
\end{practiceans}

% ---- Q58
\Needspace{9\baselineskip}
\begin{practiceans}
  \correct{B}
  The tau is a lepton, like the electron and muon, so it is fundamental.
  \begin{itemize}
    \item \xmark\ A and D: protons and neutrons are baryons (three quarks).
    \item \xmark\ C: a pion is a meson (a quark and an antiquark).
  \end{itemize}
\end{practiceans}

% ---- Q59
\Needspace{9\baselineskip}
\begin{practiceans}
  \correct{A}
  A neutron is udd: charge $+\frac{2}{3} - \frac{1}{3} - \frac{1}{3} = 0$.
  \begin{itemize}
    \item \xmark\ B: uud is a proton (charge $+1$). \quad \xmark\ C: uuu has charge $+2$.
    \item \xmark\ D: ddd has charge $-1$.
  \end{itemize}
\end{practiceans}

% ---- Q60
\Needspace{9\baselineskip}
\begin{practiceans}
  \correct{A}
  Charge $= 3 \times \left(-\frac{1}{3}e\right) = -e$.
  \begin{itemize}
    \item \xmark\ B: the charge of one strange quark only. \quad \xmark\ C and D: the charges do not add up.
  \end{itemize}
\end{practiceans}

% ---- Q61
\Needspace{9\baselineskip}
\begin{practiceans}
  \correct{B}
  A down quark changes into an up quark by emitting a $\mathrm{W}^-$ boson, which decays into an
  electron and an electron antineutrino. Charge is conserved ($0 \rightarrow +1 - 1 + 0$), and so is lepton
  number (the electron and antineutrino cancel).
  \begin{itemize}
    \item \xmark\ A: these come from beta-plus decay. \quad \xmark\ C: gluons carry the strong force
          and do not change quark type.
    \item \xmark\ D: a photon cannot carry away the charge.
  \end{itemize}
\end{practiceans}

% ---- Q62
\Needspace{9\baselineskip}
\begin{practiceans}
  \correct{D}
  Mesons (such as pions) are made of a quark and an antiquark. Baryons (such as protons) are made of
  three quarks.
  \begin{itemize}
    \item \xmark\ A: this is a baryon. \quad \xmark\ B: this is an antibaryon.
    \item \xmark\ C: two quarks do not form a particle on their own.
  \end{itemize}
\end{practiceans}

% ---- Q63
\Needspace{9\baselineskip}
\begin{practiceans}
  \correct{D}
  Strong force: gluons. Electromagnetic force: photons. Weak force: W and Z bosons.
  \begin{itemize}
    \item \xmark\ A: the two are swapped. \quad \xmark\ B and C: W and Z bosons carry the weak force.
  \end{itemize}
\end{practiceans}

% ---- Q64
\Needspace{9\baselineskip}
\begin{practiceans}
  \correct{C}
  The energy of the collision can become mass ($E = mc^2$). A particle and its antiparticle are made
  together, so charge is conserved ($0 \rightarrow +1 - 1$), and so is baryon number. This needs at least
  $2m_\mathrm{p}c^2$ (about \qty{1.9}{GeV}) of energy.
  \begin{itemize}
    \item \xmark\ A: any particles can be made, if the energy is enough and conservation laws are
          obeyed.
    \item \xmark\ B: energy can be converted into mass. \quad \xmark\ D: there must be enough energy.
  \end{itemize}
\end{practiceans}

% ---- Q65
\Needspace{9\baselineskip}
\begin{practiceans}
  \correct{C}
  $E = 2m_\mathrm{p}c^2 = 2(\num{1.673e-27})(\num{3.00e8})^2 = \qty{3.0e-10}{J}$ (about \qty{1.9}{GeV}).
  \begin{itemize}
    \item \xmark\ A: this is $2m_\mathrm{p}c$ (used $c$ instead of $c^2$).
    \item \xmark\ B: enough for one proton only. \quad \xmark\ D: doubled twice.
  \end{itemize}
\end{practiceans}

% ---- Q66
\Needspace{9\baselineskip}
\begin{practiceans}
  \correct{B}
  New particles are often very massive, so a large collision energy is needed to create them
  ($E = mc^2$). High-energy particles also have a short de Broglie wavelength ($\lambda = \frac{h}{p}$),
  so they can resolve very small structures, such as the quarks inside protons.
  \begin{itemize}
    \item \xmark\ A: nothing with mass can reach the speed of light.
    \item \xmark\ C: charge does not change with energy. \quad \xmark\ D: not related.
  \end{itemize}
\end{practiceans}

% ---- Q67
\Needspace{9\baselineskip}
\begin{practiceans}
  \correct{D}
  Before: $0$. After: $+1 + 1 + 0 = +2$. Charge is not conserved. (The real decay is
  $\mathrm{n} \rightarrow \mathrm{p} + \mathrm{e}^- + \bar{\nu}_\mathrm{e}$.)
  \begin{itemize}
    \item \xmark\ A: $+2 \rightarrow +2$. \quad \xmark\ B: $+1 \rightarrow +1 + 1 - 1 = +1$.
          \quad \xmark\ C: $+2 \rightarrow +1 + 0 + 1 = +2$.
  \end{itemize}
\end{practiceans}

% ---- Q68
\Needspace{9\baselineskip}
\begin{practiceans}
  \correct{A}
  In deep inelastic scattering experiments (Stanford, 1968), high-energy electrons were scattered by
  protons through large angles more often than expected. This showed that protons contain small,
  point-like charged particles: quarks.
  \begin{itemize}
    \item \xmark\ B, C and D: these are true, but they say nothing about the inside of the proton.
  \end{itemize}
\end{practiceans}

% ---- Q69
\Needspace{9\baselineskip}
\begin{practiceans}
  \correct{B}
  The strong force acts only on quarks (and on the particles made of them, the hadrons). The electron is a
  lepton, with no quarks, so it does not feel the strong force.
  \begin{itemize}
    \item \xmark\ A, C and D: protons and neutrons (baryons) and pions (mesons) are made of quarks.
  \end{itemize}
\end{practiceans}

% ---- Q70
\Needspace{9\baselineskip}
\begin{practiceans}
  \correct{C}
  The total energy is the rest energy of both particles, $2m_\mathrm{e}c^2$. Two photons share it
  equally (they go in opposite directions, so the total momentum stays zero):
  $E = m_\mathrm{e}c^2 = (\num{9.109e-31})(\num{3.00e8})^2 = \qty{8.2e-14}{J}$ (\qty{0.511}{MeV}).
  \begin{itemize}
    \item \xmark\ A: this is $m_\mathrm{e}c$. \quad \xmark\ B: halved the energy of one particle.
    \item \xmark\ D: this is the total energy of both photons.
  \end{itemize}
\end{practiceans}

\end{practiceanswers}
\end{document}
